\scanpage{76}
% archive sheet number: 75

bilité. Si \(p:\unsure{\mathcal E|E}\unsure{-|\longrightarrow}\unsure{B|\mathcal B}\) ne satisfait pas cette condition, une modification triviale de \(\mathcal E\) (\unsure{mod équivalence|une équivalence|aut-équivalence|avec équivalence}) la rend satisfaite.

Partant d'un foncteur fidèle transportable, on définit \deleted{sur \(B^{\mathrm{is}}\) une catégorie} un \uline{foncteur} sur \(B^{\mathrm{is}}\)

\[
\begin{aligned}
b&\unsure{\mapsto|\sim}\text{ Classes d'équiv \unsure{à|:} isom près de }\mathcal E_b\\
 &=\text{``Structures d'espèce }p\text{ sur }b\text{''}
\end{aligned}
\quad\left|\begin{gathered}\text{à valeurs dans les}\\\text{ensembles}\\\inserted{\text{ordonnés}}\\
(\text{\uline{str. \(\pm\) fine}}]\end{gathered}\right.
\]

(via ``\uline{transport de structure}'') de façon qu'un objet de \(\mathcal E\) se reconstitue à isom \inserted{unique} près comme objet de \(B\) ``muni d'une structure d'espèce \(p\)''. De façon précise, ceci permet de reconstituer la catégorie \(\mathcal E^{\unsure{\mathrm{is}|\Delta}}\) avec son foncteur \(\mathcal E^{\mathrm{is}}\xrightarrow{p^{\mathrm{is}}}B^{\mathrm{is}}\).

Pour pouvoir reconstituer \(\mathcal E\) lui-\(=\), il faudrait que \(\mathcal E\) soit fibrant (\unsure{disons|cloisonné|clésé}) sur \(B\), de façon qu'on pourra définir la notion de ``structure image inverse'' d'une structure d'espèce \(p\) sur \(c\in\operatorname{Ob} B\), par une flèche \(b\unsure{-|\longrightarrow}c\)'', de façon à avoir \inserted{\(=\)} un foncteur \(\Sigma:B^{\unsure{o|0|\circ}}\longrightarrow(\text{ens ordonnés})\) [et non \unsure{seult} \deleted{\unsure{sur}} \(B^{\mathrm{is}}\longrightarrow(\text{ens ordonnés})\)]. On reconstitue ainsi \(\mathcal E\) à \(B\)-équivalence près par la connaissance de \(\Sigma\), et \(=\) à \(B\)-isom près si \(\mathcal E\) est réduite.

Ainsi, \unsure{mod|une} \(B\)-équiv. (\(B\)-isom), les catégories \uline{fibrées} sur \(B\), à projections \uline{fidèles} \inserted{(resp. et réduites)}, s'interprètent en termes de foncteurs \inserted{\unsure{contr.}} \(B^o\longrightarrow(\unsure{\text{Ens ord}|\text{ens. ord.}})\), et pour l'exprimer le langage des structures est assez suggestif.
